% Section 3 targets. This file is included by main.tex.

\setcounter{section}{2}
\section{Problem Statement}

\begin{definition}[Allocation instance]\label{def:allocation-instance}
% lea: define
% lea: label=AllocationInstance
% lea: context={Create a structure named AllocationInstance. Bundle finite node and partition types with DecidableEq instances, capacity and size functions into Nat, request frequencies indexed by partition and node into Nat, a communication-cost function into Nat, and an exact replication parameter k : Nat. Do not build feasibility assumptions into the structure.}
An \emph{allocation instance} $I$ consists of finite types $N_I$ of storage
nodes and $P_I$ of data partitions, together with functions
\[
  \operatorname{capacity}_I : N_I \to \mathbb{N}, \qquad
  \operatorname{size}_I : P_I \to \mathbb{N},
\]
\[
  \operatorname{requests}_I : P_I \to N_I \to \mathbb{N}, \qquad
  \operatorname{communicationCost}_I : P_I \to \mathbb{N},
\]
and an exact replication parameter $k_I\in\mathbb{N}$. Equality on $N_I$ and
$P_I$ is decidable. No feasibility assumption is part of the data of an
allocation instance.
\end{definition}

\begin{definition}[Data assignment]\label{def:data-assignment}
% lea: define
% lea: label=DataAssignment
% lea: uses={AllocationInstance}
% lea: context={Define DataAssignment (I : AllocationInstance) as a function from I's partition type to Finset of I's node type. Use the primary declaration name DataAssignment.}
For an allocation instance $I$, a \emph{data assignment} is a function
\[
  D : P_I \longrightarrow \mathcal{P}_{\mathrm{fin}}(N_I).
\]
The condition $n\in D(p)$ means that node $n$ stores a copy of partition $p$.
\end{definition}

\begin{definition}[Exact $k$-safety]\label{def:exact-k-safety}
% lea: define
% lea: label=ExactKSafety
% lea: uses={AllocationInstance, DataAssignment}
% lea: context={Define a Prop-valued predicate ExactKSafety (I : AllocationInstance) (D : DataAssignment I). It should state that the card of D p equals I.k for every partition p. This is exact replication, not at-least replication.}
A data assignment $D$ for $I$ has \emph{exact $k$-safety} when every partition
is stored on exactly $k_I$ nodes:
\[
  \forall p\in P_I, \qquad |D(p)|=k_I.
\]
\end{definition}

\begin{definition}[Storage feasibility]\label{def:respects-storage}
% lea: define
% lea: label=RespectsStorage
% lea: uses={AllocationInstance, DataAssignment}
% lea: context={Define a Prop-valued predicate RespectsStorage (I : AllocationInstance) (D : DataAssignment I). Use Finset.univ filtered by membership of a node in D p, and require the sum of I.size p to be at most I.capacity n for every node n.}
A data assignment $D$ \emph{respects storage} when, for every node $n$, the
combined size of the partitions assigned to $n$ does not exceed its capacity:
\[
  \forall n\in N_I, \qquad
  \sum_{\substack{p\in P_I\\ n\in D(p)}}
    \operatorname{size}_I(p)
  \leq \operatorname{capacity}_I(n).
\]
\end{definition}

\begin{definition}[Processing cost]\label{def:processing-cost}
% lea: define
% lea: label=ProcessingCost
% lea: uses={AllocationInstance, DataAssignment}
% lea: context={Define ProcessingCost (I : AllocationInstance) (D : DataAssignment I) : Nat as a finite double sum over nodes and partitions. Include I.requests p n times I.communicationCost p exactly when n is not in D p.}
The \emph{processing cost} of an assignment $D$ is the total communication cost
of requests made at nodes that do not store the requested partition:
\[
  \operatorname{ProcessingCost}_I(D)
  = \sum_{n\in N_I}\sum_{\substack{p\in P_I\\ n\notin D(p)}}
      \operatorname{requests}_I(p,n)
      \operatorname{communicationCost}_I(p).
\]
\end{definition}

\begin{definition}[Feasible assignment]\label{def:feasible-assignment}
% lea: define
% lea: label=FeasibleAssignment
% lea: uses={ExactKSafety, RespectsStorage}
% lea: context={Define FeasibleAssignment (I : AllocationInstance) (D : DataAssignment I) : Prop as the conjunction ExactKSafety I D and RespectsStorage I D. Reuse the imported project definitions rather than restating them.}
A data assignment is \emph{feasible} when it has exact $k$-safety and respects
all storage constraints:
\[
  \operatorname{FeasibleAssignment}_I(D)
  \quad\Longleftrightarrow\quad
  \operatorname{ExactKSafety}_I(D)
  \wedge \operatorname{RespectsStorage}_I(D).
\]
\end{definition}

\begin{definition}[Optimal assignment]\label{def:optimal-assignment}
% lea: define
% lea: label=IsOptimalAssignment
% lea: uses={FeasibleAssignment, ProcessingCost}
% lea: context={Define IsOptimalAssignment (I : AllocationInstance) (D : DataAssignment I) : Prop. Require D to be feasible and its ProcessingCost to be at most that of every feasible assignment D'. Do not assert that such an assignment exists.}
A data assignment $D$ is an \emph{optimal solution} to the assignment problem
for $I$ when $D$ is feasible and no feasible assignment has lower processing
cost:
\[
  \begin{aligned}
    &\operatorname{FeasibleAssignment}_I(D) \\
    &\quad\wedge\ \forall D',\
      \operatorname{FeasibleAssignment}_I(D') \longrightarrow \\
    &\qquad \operatorname{ProcessingCost}_I(D)
      \leq \operatorname{ProcessingCost}_I(D').
  \end{aligned}
\]
This is a predicate on candidate assignments; it does not claim that a feasible
or optimal assignment exists.
\end{definition}

\begin{definition}[Quadratic unconstrained binary optimization]\label{def:qubo}
% lea: define
% lea: label=QUBO
% lea: context={Create a structure QUBO (V : Type*) parameterized by a finite linearly ordered variable type V. Store a real-valued weight function on pairs of variables. Define QUBO.energy on assignments x : V -> Fin 2 by summing w i j times x_i times x_j over i <= j, and define QUBO.IsOptimal. Keeping V as a parameter will let later files instantiate QUBO with AllocationQUBOVariable.}
For a finite linearly ordered type $V$, a \emph{quadratic unconstrained binary
optimization problem} (QUBO) consists of real weights
$w:V\to V\to\mathbb{R}$. A binary assignment $x:V\to\{0,1\}$ has energy
\[
  E_w(x)=\sum_{\substack{i,j\in V\\ i\leq j}} w(i,j)x(i)x(j).
\]
The assignment $x$ is optimal when
\[
  \forall y:V\to\{0,1\}, \qquad E_w(x)\leq E_w(y).
\]
The order restriction $i\leq j$ ensures that each unordered pair contributes
once while permitting diagonal terms $i=j$.
\end{definition}
